\documentclass[12pt, letterpaper, spanish]{article}

\usepackage{hyperref} % para manejar hipervinculos
\usepackage[spanish]{babel} % Para palabras en espanhol
\usepackage{graphicx} % LaTeX package to import graphics
\graphicspath{{images/}} % configuring the graphicx package
\usepackage[spanish]{babel} % para que LaTeX corte palabras
\usepackage[dvipsnames]{xcolor} % color en el texto

\begin{document}
% Portada
\begin{titlepage}
  \begin{center}
      \Large{Universidad Catolica "Nuestra señora de la Asuncion" \\
      Facultad de ciencias y tecnologia \\
      Complementos de Informatica}
      \includegraphics[width=0.8\textwidth]{UcaLogo.jpg}
      \LARGE{\textbf{Propuesta de Problemática y Solución}}
      \Large{Entrega 1 - Data mining}
      \vspace{1cm}
  \end{center}
      \large
      \textbf{Alumno: }Alain Vega \\
      \textbf{Profesor: }Wilfrido Felix Inchaustti Martínez
      \vfill
      \hfill{11 de septiembre de 2023}
\end{titlepage}

\newpage
\tableofcontents % indice
\newpage

% Contenidos
\section{Caso de estudio: Insuficiencia cardiaca}
\subsection{Problema}
Las enfermedades cardiovasculares (ECV) son la principal causa de muerte a nivel mundial, 
cobrando la vida de aproximadamente 17.9 millones de personas cada año, 
lo que representa el 31\% de todas las muertes en todo el mundo. 
Cuatro de cada 5 muertes por ECV son causadas por ataques cardíacos y accidentes cerebrovasculares, 
y un tercio de estas muertes ocurren prematuramente en personas menores de 70 años. 

La insuficiencia cardíaca es un evento común causado por las ECV.
Las personas con enfermedades cardiovasculares o que tienen un alto riesgo cardiovascular 
(debido a la presencia de uno o más factores de riesgo como hipertensión, diabetes, 
hiperlipidemia o enfermedades ya establecidas) necesitan una detección temprana y gestión, 
en la cual un modelo de aprendizaje automático puede ser de gran ayuda.
\section{Determinar objetivos del negocio}
\subsection{Background}
La insuficiencia cardíaca se produce cuando el músculo cardíaco no bombea sangre 
de la manera que debería. Cuando esto sucede, la sangre a menudo retrocede y 
el líquido se puede acumular en los pulmones, lo que causa falta de aliento.

El conjunto de datos a considerar, analiza 11 caracteristicas de un total de 918 personas.
\subsection{Objetivos del negocio}
  \begin{itemize}
    \item{\textbf{Reducción de costos de atención médica: }
    al prevenir diagnósticos fallidos y tratamientos innecesarios, se reducen los costos
    asociados con el tratamiento de insuficiencia cardíaca y enfermedades relacionadas.}
    \item{\textbf{Mejorar la atención médica:} 
    al identificar de manera más efectiva a aquellos en riesgo de insuficiencia cardíaca. 
    Esto conlleva a un tratamiento más oportuno y eficaz.}
    \item{\textbf{Aumentar la calidad de vida de los pacientes:} 
    al permitirles recibir atención preventiva y seguimiento temprano gracias a la prediccion.}
  \end{itemize}
\subsection{Criterios de exito del negocio}
  \begin{itemize}
    \item{\textbf{Reduccion de costos: }dismunir los costos de tratamientos medicos un 20\%
    para los pacientes}
    \item{\textbf{Impacto en la vida de los pacientes: }afectar positivamente la calidad de vida
    de los pacientes.}
  \end{itemize}

\section{Valoracion de la situacion}
\subsection{Inventario de recursos}
Se cuenta con los siguientes datos medicos: \emph{edad, sexo, tipo de dolor en el pecho, 
presión arterial en reposo, colesterol, azúcar en sangre en ayunas, 
resultados del electrocardiograma en reposo, frecuencia cardíaca máxima alcanzada, 
angina inducida por el ejercicio, pico antiguo = ST, 
pendiente del segmento ST durante un ejercicio físico máximo en una prueba de esfuerzo cardíaco, 
sufrio de insuficiencia cardiaca.}

Que serviran para el entrenamiento y pruebas del modelo.
\subsection{Requisitos, supuestos y restricciones}
\begin{itemize}
  \item{\textbf{Requisitos: }deteccion temprana de posible insuficiencia cardiaca
  con una tasa de falsos negativos cada vez menor en proporcion a la
  cantidad de pacientes analizados.} 
  \item{\textbf{Supuestos: }los datos iniciales proveidos son de fiar, representan
  las caracteristicas principales a analizar.}
  \item{\textbf{Restricciones: }la fiabilidad del modelo nunca debe bajar del 90\%}
\end{itemize}
\subsection{Riesgos y Contingencias}
\begin{itemize}
  \item{\textbf{Riesgos: }algunas afecciones cardíacas progresivamente dejan 
  el corazón demasiado débil o rígido como para llenarse y bombear sangre de forma apropiada. 
  Por lo que los síntomas pueden desarrollarse lentamente. Sin embargo a veces, 
  los síntomas de la insuficiencia cardíaca comienzan repentinamente.
  
  Entonces existe el riesgo de que el conjunto de datos a analizar no contemple todas
  las variables para un analisis exacto.} 
  \item{\textbf{Contingencias: }agregar nuevas caracteristicas relevantes al conjunto de datos,
  ademas de incluir mas registros con el fin de tener una muestra mas grande y variada.}
\end{itemize}
\subsection{Costes y Beneficios}
De acuerdo a un estudio realizado en Estados Unidos, el costo promedio por año de pacientes
con insuficiencia cardiaca en dicho pais es de \textcolor{red}{\$18.953} de los cuales el 
42,8\% (\$8114) son costos hospitalarios.

Gracias a una deteccion temprana es posible evitar que el paciente se hospitalice por 
un cuadro avanzado de esta afeccion. Por lo que es posible reducir el costo anual a 
\textbf{\$18.953 * (1 - 0,428) = \textcolor{YellowGreen}{\$10.841,116}}.

Esto es una reduccion del 57,2\% de los costos promedios anuales en el tratamiento.

\section{Determinar los objetivos del DM}
\subsection{Metas de Data Mining}
\begin{itemize}
  \item{\textbf{Precision del modelo: }obtener un modelo fiable, que posea una precision 
  mayor o igual al 95\% }
  \item{\textbf{Implementación Clínica: }implementar el modelo en mas de un entorno clínico real 
  para ayudar en la detección temprana y la gestión de la insuficiencia cardíaca.}
\end{itemize}
\subsection{Criterios de éxito de DM}
Se establecen como criterios de exito, los siguientes:
\begin{itemize}
  \item{\textbf{Precision del modelo: }obtener un modelo fiable, que posea una precision de
  por lo menos el 90\% }
  \item{\textbf{Reduccion de costos medicos: }dismunir los costos de tratamientos medicos 
  un 57,2\% para los pacientes}
\end{itemize}

\section{Bibliografia}
\begin{itemize}
    \item \href{https://www.sngular.com/es/crisp-dm-fase-i-comprension-del-negocio/?authuser=0}
    {Metodologia CRISP-DM}
    \item \href{https://www.kaggle.com/datasets/fedesoriano/heart-failure-prediction}
    {Conjunto de datos de insuficiencia cardiaca - \emph{kaggle}}
    \item \href{https://www.ncbi.nlm.nih.gov/pmc/articles/PMC8730466/}{Costo promedio anual de la insuficiencia cardiaca}
\end{itemize}

\end{document}